%%%%%%%%%%%%%%%%%%%%%%%%%%%%%%%%%%%%%%%%%%%%%%%%%%%%%%%%%%%%%%%%%%%%%%%%%%%%%%%
%% Lista de Símbolos (pacote nomencl)
%%
%% Contém os símbolos que o texto define e usa. O prefixo numérico mantém a ordem de
%% primeira ocorrência, exigida pela NBR 14724; ao acrescentar um símbolo, escolha o
%% prefixo pela posição dele no texto. A macro \nomunit acrescenta a unidade alinhada à
%% direita.
%%%%%%%%%%%%%%%%%%%%%%%%%%%%%%%%%%%%%%%%%%%%%%%%%%%%%%%%%%%%%%%%%%%%%%%%%%%%%%%

\renewcommand{\nomname}{Lista de símbolos}

%----------------------------------------------
\newcommand{\nomunit}[1]{%
\renewcommand{\nomentryend}{\hspace*{\fill}#1}}
%----------------------------------------------

%% --- Vibração e descritores --------------------------------------------------
\nomenclature[001]{\(x(t)\)}{Deslocamento\nomunit{\si{\meter}}}
\nomenclature[002]{\(t\)}{Tempo\nomunit{\si{\second}}}
\nomenclature[003]{\(v(t)\)}{Velocidade\nomunit{\si{\meter\per\second}}}
\nomenclature[004]{\(a(t)\)}{Aceleração\nomunit{\si{\meter\per\second\squared}}}
\nomenclature[005]{\(f\)}{Frequência\nomunit{\si{\hertz}}}
\nomenclature[006]{\(X\)}{Amplitude de deslocamento\nomunit{\si{\meter}}}
\nomenclature[007]{\(\omega\)}{Frequência angular\nomunit{\si{\radian\per\second}}}
\nomenclature[008]{\(N\)}{Número de amostras na janela}
\nomenclature[009]{\(x_n\)}{Amostra de índice \(n\) da janela}
\nomenclature[010]{\(\bar{x}\)}{Média do sinal na janela}
\nomenclature[011]{\(\sigma\)}{Desvio-padrão do sinal na janela\nomunit{\si{\meter\per\second\squared}}}
\nomenclature[012]{\(x_{\mathrm{rms}}\)}{Valor eficaz do sinal na janela\nomunit{\si{\meter\per\second\squared}}}
\nomenclature[013]{\(\kappa\)}{Curtose do sinal na janela}
\nomenclature[014]{\(x_{\mathrm{pico}}\)}{Valor de pico do sinal na janela}

%% --- Rolamento e aquisição ---------------------------------------------------
\nomenclature[022]{\(f_r\)}{Frequência de rotação do eixo\nomunit{\si{\hertz}}}
\nomenclature[023]{\(N_b\)}{Número de esferas do rolamento}
\nomenclature[024]{\(d\)}{Diâmetro da esfera\nomunit{\si{\milli\meter}}}
\nomenclature[025]{\(D_p\)}{Diâmetro primitivo do rolamento\nomunit{\si{\milli\meter}}}
\nomenclature[026]{\(\theta\)}{Ângulo de contato do rolamento\nomunit{\si{\radian}}}
\nomenclature[027]{\(f_s\)}{Taxa de amostragem\nomunit{\si{\hertz}}}
\nomenclature[028]{\(f_{\mathrm{ap}}\)}{Frequência aparente sob subamostragem sem filtro\nomunit{\si{\hertz}}}
\nomenclature[029]{\(X_k\)}{Coeficiente de ordem \(k\) da DFT}
\nomenclature[030]{\(\Delta f\)}{Resolução em frequência da DFT\nomunit{\si{\hertz}}}
\nomenclature[031]{\(B\)}{Largura da banda de ruído\nomunit{\si{\hertz}}}

%% --- Detector e avaliação ----------------------------------------------------
\nomenclature[032]{\(h(x)\)}{Comprimento do caminho de isolamento de \(x\)}
\nomenclature[033]{\(\psi\)}{Tamanho da subamostra de cada árvore}
\nomenclature[034]{\(c(\psi)\)}{Comprimento médio de caminho para subamostra de \(\psi\) pontos}
\nomenclature[035]{\(s(x,\psi)\)}{Escore de anomalia do Isolation Forest}
\nomenclature[036]{\(\alpha\)}{Fração da operação normal acima do limiar}
\nomenclature[037]{\(F_1\)}{Medida que combina precisão e TPR}
\nomenclature[038]{\(\beta\)}{Limite superior de falso alarme da área parcial}
\nomenclature[0385]{\(A_\beta\)}{Área sob a curva ROC entre 0 e \(\beta\)}
\nomenclature[039]{\(\mathrm{pAUC}(\beta)\)}{Área parcial padronizada sob a curva ROC}

%% --- Enlace LoRa -------------------------------------------------------------
\nomenclature[040]{\(R_s\)}{Taxa de símbolos LoRa\nomunit{\si{\per\second}}}
\nomenclature[041]{\(T_s\)}{Duração de um símbolo LoRa\nomunit{\si{\second}}}
\nomenclature[042]{\(T_{\mathrm{pacote}}\)}{Duração de um pacote LoRa\nomunit{\si{\second}}}
\nomenclature[043]{\(n_{\mathrm{pre}}\)}{Número de símbolos do preâmbulo}

%% --- Modelo de sinal (seção de métodos) --------------------------------------
\nomenclature[047]{\(A_k\)}{Amplitude da \(k\)-ésima harmônica da rotação\nomunit{\si{\meter\per\second\squared}}}
\nomenclature[048]{\(\varphi_k\)}{Fase da \(k\)-ésima harmônica\nomunit{\si{\radian}}}
\nomenclature[049]{\(D\)}{Amplitude do impacto do defeito\nomunit{\si{\meter\per\second\squared}}}
\nomenclature[050]{\(\zeta_d\)}{Taxa de amortecimento do impacto\nomunit{\si{\per\second}}}
\nomenclature[051]{\(f_{\mathrm{res}}\)}{Frequência da ressonância estrutural excitada\nomunit{\si{\hertz}}}
\nomenclature[052]{\(t_i\)}{Instante do \(i\)-ésimo impacto\nomunit{\si{\second}}}
\nomenclature[053]{\(\varepsilon(t)\)}{Ruído gaussiano branco\nomunit{\si{\meter\per\second\squared}}}
\nomenclature[054]{\(u(\cdot)\)}{Degrau unitário}
\nomenclature[055]{\(\delta_i\)}{Desvio aleatório que modela o escorregamento das esferas\nomunit{\si{\second}}}
